\documentclass[12pt]{article}
\usepackage[a4paper,margin=1in]{geometry}
\usepackage{graphicx}
\usepackage[export]{adjustbox}
\usepackage{hyperref}
\usepackage{fontspec}
\setmainfont{DejaVu Sans}
\usepackage{amsmath}
\usepackage{amssymb}
\usepackage{listings}

\begin{document}
APEX INSTITUTE OF TECHNOLOGY — SCHOOL OF COMPUTING
Master of Science in Data Science \& Artificial
Intelligence (MS-DSAI)
Academic Year: 2026–
2027
•
Degree Level: Postgraduate (Master of
Science)
•
Total Credits: 36 US Credits / 90
ECTS
1. Executive Summary \& Program Vision
Program Mission: The Master of Science in Data Science \& Artificial Intelligence (MS-DSAI) at Apex Institute of
Technology is an advanced professional degree program designed to bridge the gap between theoretical
computing foundation and high-impact industrial deployment. The curriculum empowers graduates with
rigorous mathematical foundations, scalable machine learning systems engineering, and ethical AI governance
capabilities required to lead complex data-driven transformations in modern global enterprises.
36
TOTAL CREDITS
2 Years
FULL-TIME DURATION
98\%
PLACEMENT RATE
1:8
FACULTY RATIO
2. Program Educational Objectives (PEOs)
Upon successful completion of the MS-DSAI program, graduates are expected to achieve the following career and
professional milestones within 3 to 5 years of graduation:
PEO 1 (Technical Mastery): Lead the architecture, design, and deployment of complex artificial intelligence systems and
enterprise data platforms handling multi-terabyte computational pipelines.
PEO 2 (Applied Innovation): Formulate original algorithmic solutions to ambiguous, real-world domain problems across
sectors such as healthcare, finance, autonomous systems, and modern logistics.
PEO 3 (Ethical Leadership): Establish robust AI governance, algorithmic transparency, privacy preservation, and ethical
compliance frameworks within technological organizations.
PEO 4 (Collaborative Excellence): Communicate complex quantitative and computational insights clearly to non-
technical executive stakeholders, cross-functional engineering teams, and regulatory bodies.
3. Core Pillars of the Curriculum
1. Rigorous Mathematics
Deep coverage of vector calculus,
high-dimensional probability, convex
optimization, stochastic processes,
and linear algebra underpinning
modern deep architectures.
2. Systems Engineering
Practical experience with distributed
computing, real-time data streaming,
GPU acceleration clusters, MLOps,
containerization, and microservices.
3. Responsible AI
Direct integration of algorithmic
fairness, explainability (XAI),
federated learning, data privacy laws,
and regulatory safety compliance
throughout all tracks.
Master of Science in Data Science \& Artificial Intelligence — Academic Handbook
Page 1 of 12

4. Curriculum Architecture \& Credit Distribution
The MS-DSAI degree requires the completion of 36 graduate credits. Students progress through a structured
sequence consisting of Core Foundations (18 credits), Specialized Track Electives (9 credits), Industry Practicum /
Graduate Seminar (3 credits), and a Culminating Master's Thesis or Capstone Project (6 credits).
Category
Required
Credits
ECTS
Equivalent
Description
Core Foundations
18 Credits
45 ECTS
6 mandatory core courses covering math, algorithmics, machine
learning, deep learning, distributed systems, and NLP/CV.
Specialization
Electives
9 Credits
22.5 ECTS
3 depth courses selected from one designated specialization track
(ML Systems, Cloud Data Architecture, or AI Governance).
Graduate Seminar /
Practicum
3 Credits
7.5 ECTS
Industry mentorship series, research methodology, and emerging
topic seminars.
Thesis / Capstone
Project
6 Credits
15 ECTS
Two-semester original research thesis or enterprise client capstone
project.
Total Degree
Requirement
36 Credits
90 ECTS
4 Semesters (24 Months Full-Time / 36 Months Part-Time)
5. Academic Progression Roadmap
The diagram below depicts the recommended sequence of study for full-time graduate students across the four
academic semesters:
SEMESTER 1 (10 Cr)
SEMESTER 2 (10 Cr)
SEMESTER 3 (9 Cr)
SEMESTER 4 (7 Cr)
DSAI-501 (3 Cr)
Applied Math \& Stats
DSAI-502 (3 Cr)
Data Structures \& Python
DSAI-503 (3 Cr)
Machine Learning Foundations
DSAI-500 (1 Cr) Colloquium
DSAI-504 (3 Cr)
Deep Learning \& Architectures
DSAI-505 (3 Cr)
Distributed Big Data Systems
DSAI-506 (3 Cr)
Computer Vision \& NLP
DSAI-590 (1 Cr) Ethics Lab
Elective Track I (3 Cr)
Specialization Deep Dive
Elective Track II (3 Cr)
Specialization Deep Dive
DSAI-698 (3 Cr)
Capstone / Thesis Phase I
Elective Track III (3 Cr)
Specialization Deep Dive
DSAI-699 (3 Cr)
Thesis / Capstone II
Final Defense \& Oral Exam
DSAI-600 (1 Cr) Portfolio
6. Key Graduation Requirements Summary
Minimum Cumulative GPA: Students must maintain a minimum cumulative Grade Point Average (GPA) of 3.00 out of
4.00 across all attempted graduate coursework.
Core Competency Standard: A grade of 'B' (3.0) or higher is required in all six Core Foundation courses.
Residency Requirement: At least 27 credits must be completed in physical or official synchronous residence at Apex
Institute of Technology.
Thesis / Capstone Defense: Successful public defense and committee sign-off on the Master's Thesis or Enterprise
Capstone Project.
Master of Science in Data Science \& Artificial Intelligence — Academic Handbook
Page 2 of 12

7. Core Course Syllabi — Year 1, Semester 1
The first semester establishes the foundational mathematical principles and high-performance computing
competencies required for advanced machine learning research and applications.
DSAI-501 — Applied Probability \& Mathematical Statistics for Data Science
3 Credits • Core
Course Overview: Advanced treatment of multivariate probability distribution theory, maximum likelihood
estimation, Bayesian inference, Markov Chains, expectation-maximization (EM) algorithm, hypothesis testing,
and high-dimensional vector spaces.
Learning Outcomes: Formulate rigorous probabilistic models; implement Monte Carlo sampling methods; derive
closed-form loss functions for parameter estimation.
Probability Theory Bayesian Inference Markov Chains Vector Calculus
Assessment: Homework Assignments (25\%), Midterm Exam (35\%), Final Theoretical Examination (40\%).
DSAI-502 — Advanced Algorithmic Data Structures \& Python Systems
3 Credits • Core
Course Overview: Intensive systems-level programming in Python and C++ for data engineering. Covers
asymptotic memory analysis, vectorized array computations, graph processing, parallel processing, and object-
oriented software design patterns.
Learning Outcomes: Build memory-efficient algorithmic data structures; optimize memory footprints for matrix
transformations; write parallel C-extensions for Python environments.
Python C-API NumPy Internal Mechanics Graph Algorithms Memory Profiling
Assessment: Bi-weekly Coding Labs (40\%), System Architecture Project (30\%), Final Exam (30\%).
DSAI-503 — Foundations of Machine Learning \& Statistical Learning
3 Credits • Core
Course Overview: Comprehensive mathematical analysis of supervised and unsupervised learning algorithms.
Topics include linear regression, logistic regression, SVMs, decision tree ensembles (Random Forests, XGBoost),
dimensionality reduction (PCA, t-SNE), and regularization techniques (L1/L2, ElasticNet).
Learning Outcomes: Implement core ML algorithms from mathematical principles without reliance on high-level
libraries; evaluate bias-variance trade-offs mathematically.
Supervised Learning Convex Optimization Ensemble Methods Statistical Bounds
Assessment: Algorithm Construction Projects (35\%), Midterm Exam (30\%), Final Project \& Paper (35\%).
Master of Science in Data Science \& Artificial Intelligence — Academic Handbook
Page 3 of 12

8. Core Course Syllabi — Year 1, Semester 2
The second semester advances into modern artificial intelligence architectures, distributed big data infrastructure,
and complex multi-modal information processing.
DSAI-504 — Deep Learning Architecture \& Neural Networks
3 Credits • Core
Course Overview: Deep study of artificial neural networks, automatic differentiation engines, convolutional
networks (CNNs), sequence models (RNNs, LSTMs), and self-attention Transformer mechanisms. Hands-on
model development using PyTorch.
Learning Outcomes: Design custom neural architectures; implement backpropagation from scratch; optimize
hyperparameter search strategies on distributed GPU clusters.
PyTorch Transformers Attention Mechanisms Gradient Descent
Assessment: Deep Learning Benchmarking Labs (30\%), Model Architecture Challenge (30\%), Final Exam (40\%).
DSAI-505 — Distributed Big Data Systems \& Cloud Computing
3 Credits • Core
Course Overview: Engineering scalable analytics pipelines across cloud and edge clusters. Deep exploration of
Apache Spark, Hadoop MapReduce, distributed file systems (HDFS/S3), distributed SQL databases, and stream
processing engines (Kafka, Flink).
Learning Outcomes: Architect fault-tolerant distributed data processing jobs; optimize execution DAGs; query
multi-terabyte datasets under sub-second latency constraints.
Apache Spark Kafka Streaming Cloud Architecture Distributed SQL
Assessment: Distributed Cluster Engineering Labs (40\%), Scalability Case Study (25\%), Final Project (35\%).
DSAI-506 — Applied Computer Vision \& Natural Language Processing
3 Credits • Core
Course Overview: Dual focus on visual and linguistic intelligence. Covers object detection, semantic
segmentation, modern Large Language Models (LLMs), tokenization, embedding spaces, retrieval-augmented
generation (RAG), and fine-tuning methodologies.
Learning Outcomes: Fine-tune pre-trained vision-language models; deploy RAG pipelines over enterprise
document stores; evaluate semantic alignment metrics.
LLMs \& RAG Object Detection Vector Databases Hugging Face
Assessment: Multi-modal AI Hackathon (35\%), Research Paper Replication (25\%), Final Examination (40\%).
Master of Science in Data Science \& Artificial Intelligence — Academic Handbook
Page 4 of 12

9. Specialization Track 1: Machine Learning Systems Engineering
Track Focus: Designed for students aspiring to become Principal MLOps Engineers, AI Systems Architects, and
Machine Learning Platform Leads. Focuses on the production lifecycle, hardware acceleration, scalable model
deployment, continuous evaluation, and infrastructure monitoring.
DSAI-610 — Scalable MLOps \& Production Infrastructure
3 Credits • Track Elective
Continuous Integration / Continuous Deployment (CI/CD) for Machine Learning, feature stores (Feast), model
registries (MLflow), Kubernetes model serving, dynamic auto-scaling, model drift detection, and automated
rollback strategies.
DSAI-611 — Reinforcement Learning \& Autonomous Agents
3 Credits • Track Elective
Markov Decision Processes (MDPs), Q-learning, Policy Gradient methods, Deep Q-Networks (DQN), Proximal
Policy Optimization (PPO), multi-agent reinforcement learning, and RLHF (Reinforcement Learning from Human
Feedback) for alignment.
DSAI-612 — Generative AI \& Large Language Model Engineering
3 Credits • Track Elective
Architectural innovations in Generative Adversarial Networks (GANs), Diffusion Models, Transformer variants,
parameter-efficient fine-tuning (LoRA, QLoRA), quantizations (AWQ, GGUF), and high-throughput inference
engines (vLLM, TensorRT-LLM).
Track 1 Computing Hardware \& Infrastructure Stack
Students enrolled in Track 1 receive dedicated allocation compute tokens on the Apex AI GPU Supercomputing
Cluster (8x NVIDIA H100 nodes per pod), enabling hands-on experience with multi-node distributed training.
Orchestration \& Serving Stack
Kubernetes \& Kubeflow Pipelines
Triton Inference Server
Ray Cluster Management
Optimization \& Storage
ONNX Runtime \& TensorRT
Pinecone \& Milvus Vector DBs
Weave \& Weights \& Biases
Master of Science in Data Science \& Artificial Intelligence — Academic Handbook
Page 5 of 12

10. Specialization Track 2: Data Systems \& Cloud Analytics
Track Focus: Engineered for future Data Platform Architects, Chief Data Officers, and Cloud Analytics Leads.
Focuses on high-throughput data pipelines, enterprise data warehousing, real-time event streaming
architectures, and modern data lakehouse designs.
DSAI-620 — High-Performance Database Systems \& NoSQL
3 Credits • Track Elective
Internal storage engines (B-Trees, LSM-Trees), WAL recovery, ACID transaction isolation, distributed consensus
protocols (Raft, Paxos), key-value stores, document databases (MongoDB), wide-column stores (Cassandra), and
graph databases (Neo4j).
DSAI-621 — Real-Time Streaming Systems \& Event-Driven Architecture
3 Credits • Track Elective
Stream-table duality, event-driven microservices, Apache Flink stateful stream processing, exactly-once
processing semantics, out-of-order event handling with watermarks, and high-velocity fraud detection event
architectures.
DSAI-622 — Data Warehousing, Lakehouses \& Enterprise BI
3 Credits • Track Elective
Modern data stack architecture using Snowflake, Databricks Delta Lake, and Apache Iceberg. Dimensional
modeling, Star/Snowflake schemas, dbt transformation frameworks, semantic layers, and automated data
quality testing.
Industry System Case Studies Covered
Platform Industry
Architecture Highlight
Key Technical Objective
Financial Services
Sub-millisecond Market Data Streams
Zero-data-loss event streaming with Kafka \& Flink
E-Commerce Scale
Petabyte Data Lakehouse
ACID transactions over object storage with Apache Iceberg
Global Logistics
Real-time IoT Telemetry Engine
Time-series database partitioning \& dynamic querying
Master of Science in Data Science \& Artificial Intelligence — Academic Handbook
Page 6 of 12

11. Specialization Track 3: AI Governance, Ethics, \& Responsible Tech
Track Focus: Prepares technical leaders, regulatory consultants, AI auditors, and policy specialists. Focuses on
algorithmic fairness auditing, regulatory compliance (EU AI Act, NIST AI RMF), explainable AI (XAI), and privacy-
preserving machine learning.
DSAI-630 — AI Policy, Regulatory Frameworks \& Algorithmic Safety
3 Credits • Track Elective
In-depth analysis of global legislation including the EU AI Act, US NIST AI Risk Management Framework, and
global data privacy standards (GDPR, CCPA). Risk classification, technical auditing protocols, compliance
reporting, and liability.
DSAI-631 — Privacy-Preserving ML \& Federated Systems
3 Credits • Track Elective
Mathematical principles of Differential Privacy (DP), Homomorphic Encryption (HE), Secure Multi-Party
Computation (SMPC), federated learning aggregation algorithms (FedAvg, FedProx), and defense against
membership inference attacks.
DSAI-632 — Explainable AI (XAI) \& Algorithmic Fairness
3 Credits • Track Elective
Feature attribution mechanisms (SHAP, LIME, Integrated Gradients), counterfactual explanations, mathematical
definitions of algorithmic fairness (demographic parity, equalized odds), and bias mitigation in training pipelines.
Research Seminar Series \& Industry Guest Lectures
Track 3 features mandatory bi-weekly research colloquiums with leading ethicists, legal scholars, and senior
compliance officers from international regulatory bodies and tech companies.
Colloquium Series A: Regulatory AI
Focuses on operationalizing the EU AI Act technical
documentation requirements in high-risk deployments
(healthcare, credit scoring, hiring).
Colloquium Series B: Technical Auditing
Practical workshops on implementing automated fairness
testing suites in continuous integration pipelines using
open-source tools (AIF360, Fairlearn).
Master of Science in Data Science \& Artificial Intelligence — Academic Handbook
Page 7 of 12

12. Master's Thesis \& Capstone Project Framework
As the culminating requirement for the MS-DSAI degree, candidates must complete either an original academic
Research Thesis Track or an industry-sponsored Enterprise Capstone Project Track across Semesters 3 and 4 (6
credits total).
Option A: Academic Research Thesis
Target Audience: Students planning to pursue Ph.D.
studies or research lab roles.
Original theoretical or algorithmic research contribution.
Faculty mentor supervision + 2 committee members.
Requirement to submit a peer-reviewed paper to a top-tier
conference (NeurIPS, ICML, KDD, CVPR, ACL).
Option B: Enterprise Capstone Project
Target Audience: Students pursuing immediate industry
leadership paths.
Sponsored project solving a live enterprise industry
challenge.
Co-supervision by an industry principal engineer \& faculty
advisor.
Production-ready code repository delivery + executive
defense.
13. Evaluation Rubric \& Milestone Timeline
Milestone Phase
Target Schedule
Key Deliverables
Weight
Proposal Defense
End of Semester 3
(Week 14)
Formal Literature Review, Problem Statement, Methodology
Proposal \& Advisory Committee Approval.
15\%
Mid-Term Progress
Review
Semester 4 (Week 7)
Working prototype execution, preliminary benchmark evaluations,
code architecture review.
25\%
Final Thesis/Project
Report
Semester 4 (Week
13)
Publication-quality research manuscript or comprehensive
technical documentation book.
35\%
Public Oral Defense
Semester 4 (Week
15)
Public presentation (30 mins) + rigorous oral examination by the
graduate faculty evaluation committee.
25\%
Project Timeline Visual Diagram
1
Proposal
Sem 3, Wk 14
2
Alpha Build
Sem 4, Wk 7
3
Draft Submission
Sem 4, Wk 13
4
Final Oral Defense
Sem 4, Wk 15
Master of Science in Data Science \& Artificial Intelligence — Academic Handbook
Page 8 of 12

14. Academic Grading Scale \& Evaluation Standards
Apex Institute of Technology enforces strict academic evaluation criteria to maintain global accredited standards.
Graduate performance is graded according to the letter-grade system below:
Grade
Grade
Points
Percentage
Equivalent
Performance Definition
A
4.00
93\% – 100\%
Exemplary Mastery. Demonstrates superior theoretical understanding and
flawless practical execution.
A-
3.70
90\% – 92\%
Excellent Performance. Outstanding comprehension with minor non-critical
technical errors.
B+
3.30
87\% – 89\%
Very Good. Solid technical execution above average graduate standards.
B
3.00
83\% – 86\%
Good / Satisfactory. Meets standard requirements for core course progression.
B-
2.70
80\% – 82\%
Below Graduate Average. Credited, but triggers Academic Progress review if
cumulative GPA falls < 3.00.
C+ /
C
2.30 / 2.00
73\% – 79\%
Unsatisfactory Performance. Course must be repeated if listed as a mandatory
core course.
F
0.00
Below 70\%
Failing Grade. Zero credit awarded. Mandatory repeat required.
15. Academic Integrity \& Responsible Generative AI Policy
Strict Honor Code Policy: Academic dishonesty, including uncredited plagiarism, code fabrication, contract
cheating, and unauthorized examination assistance, results in immediate referral to the Graduate Disciplinary
Board with potential academic expulsion.
Policy on Generative AI Usage in Coursework
Permitted Use (with Citation): Utilizing LLM assistants for syntax debugging, spellchecking, and brainstorming is
permitted ONLY when explicitly authorized in writing by the course instructor. All AI-assisted code segments or text must
be explicitly disclosed in a dedicated "Generative AI Attribution Index" appended to the submission.
Prohibited Use: Generating complete assignment codebases, submitting synthetic research literature reviews, or
employing AI during closed-book theoretical examinations constitutes severe academic misconduct.
16. Academic Probation \& Remediation Protocols
Any student whose cumulative GPA falls below 3.00 at the conclusion of any semester is automatically placed on
Academic Probation. The student is assigned an Academic Advisory Specialist and granted one semester to raise
their cumulative GPA to 3.00 or higher. Failure to satisfy probation terms leads to formal academic dismissal.
Master of Science in Data Science \& Artificial Intelligence — Academic Handbook
Page 9 of 12

17. Distinguished Program Faculty
The MS-DSAI program is delivered by internationally recognized scholars, researchers, and former principal industry
scientists:
Prof. Elena Rostova, Ph.D.
Department Chair \& Professor of AI
Ph.D. Computer Science (Stanford)
Specialization: High-dimensional optimization, deep generative
architectures, and federated learning protocols. Former
Research Lead at Google Brain.
Prof. Marcus Vance, Ph.D.
Director, Center for Autonomous Systems
Ph.D. Robotics (MIT)
Specialization: Computer vision, real-time spatial sensing, deep
reinforcement learning for robotics, and autonomous agent
navigation.
Dr. Aisha Al-Hassan, Ph.D.
Associate Professor of Data Systems
Ph.D. Computer Engineering (UC Berkeley)
Specialization: Distributed database systems, stream
processing, and scalable cloud analytics architectures. Co-
inventor of key stream indexing algorithms.
Dr. Jonathan Thorne, J.D., Ph.D.
Director, Responsible AI \& Policy Lab
Ph.D. CS (Oxford), J.D. (Harvard)
Specialization: Algorithmic fairness, explainable AI, technical
auditing frameworks, and international regulatory compliance
law.
18. Associated Research Centers \& Facilities
Apex Supercomputing Core
State-of-the-art cluster hosting 64x
NVIDIA H100 GPU nodes with 100Gbps
InfiniBand interconnects, dedicated to
student research and large-scale model
pre-training.
Responsible AI Testing Lab
Interdisciplinary laboratory focused on
developing automated bias detection
toolkits, adversarial robustness
benchmarks, and privacy-preserving
federated networks.
Autonomous Vision Facility
Advanced robotics and spatial
computing arena equipped with multi-
camera motion capture rigs, LiDAR
systems, and edge GPU simulation
environments.
Active External Research Grants Showcase
NSF Frontier Grant (\#2025-AI-889): "Scalable Formal Verification for Deep Neural Networks in Critical Healthcare
Operations." (\$3.2M Award).
DARPA Autonomous Systems Initiative: "Multi-Agent Reinforcement Learning under Communication-Constrained
Environments." (\$2.8M Award).
Master of Science in Data Science \& Artificial Intelligence — Academic Handbook
Page 10 of 12


19. Tiêu chuẩn Tuyển sinh \& Quy trình Nộp hồ sơ
Chương trình MS-DSAI tìm kiếm các ứng viên có năng lực định lượng xuất sắc và nền tảng lập trình vững chắc. Việc xét tuyển được thực hiện theo hình thức cuốn chiếu (rolling basis) cho cả hai kỳ nhập học Mùa Thu và Mùa Xuân.
Yêu cầu Tiền đề Học thuật
Bằng Cử nhân Khoa học Máy tính (CS), Khoa học Dữ liệu, Toán học, Thống kê,
Vật lý, hoặc Kỹ thuật.
GPA Đại học Tối thiểu: 3.20 trên thang điểm 4.00.
Các môn học tiên quyết: Giải tích III, Đại số Tuyến tính,
Lý thuyết Xác suất, và Cấu trúc Dữ liệu bằng C++/Python.
Danh mục Hồ sơ Ứng tuyển
Bảng điểm Đại học Chính thức.
3 Thư Giới thiệu (Học thuật/Chuyên môn).
Bài luận Mục đích Học tập (1.000 từ).
Bài thi GRE General (Không bắt buộc, khuyến nghị
điểm Định lượng > 165).
TOEFL iBT (Tối thiểu 100) hoặc IELTS (Tối thiểu 7.5) đối với người không phải
là người bản ngữ tiếng Anh.
20. Học phí, Lệ phí \& Hỗ trợ Tài chính
Khoản chi phí
Chi phí trên mỗi đơn vị
Tổng chi phí toàn chương trình (36 Tín chỉ)
Học phí Sau đại học
\$1.350 mỗi tín chỉ
\$48.600
Phí Siêu máy tính \& Hạ tầng Phòng thí nghiệm
\$450 mỗi học kỳ
\$1.800
Phí Dịch vụ \& Sức khỏe Sinh viên Sau đại học
\$350 mỗi học kỳ
\$1.400
Tổng Học phí \& Lệ phí Ước tính
—
\$51.800 USD
Hỗ trợ Tài chính Sau đại học \& Học bổng Khuyến học
Trợ lý Nghiên cứu (RA)
Đài thọ 100\% học phí + trợ cấp hàng tháng \$2.800.
Được xét duyệt cạnh tranh dựa trên sự phù hợp về hướng nghiên cứu với
các dự án tài trợ của giảng viên. Yêu cầu 20 giờ/tuần làm việc tại phòng thí nghiệm.
Trợ giảng (TA)
Đài thọ từ 50\% đến 100\% học phí + trợ cấp \$2.400/tháng.
Các TA hỗ trợ giảng viên trong việc chấm điểm, hướng dẫn các buổi
thực hành phòng lab, và chủ trì các giờ tư vấn sinh viên.
Thạc sĩ Khoa học Dữ liệu \& Trí tuệ Nhân tạo — Sổ tay Học thuật
Trang 11 / 12

21. Định hướng Nghề nghiệp \& Việc làm Sau Đại học
Sinh viên tốt nghiệp chương trình MS-DSAI được tuyển dụng bởi các công ty công nghệ hàng đầu thế giới, các phòng thí nghiệm nghiên cứu, các tổ chức tài chính, và các công ty khởi nghiệp sáng tạo có tốc độ phát triển cao.
\$148k
MỨC LƯƠNG KHỞI ĐIỂM TRUNG BÌNH
98\%
CÓ VIỆC LÀM TRONG VÒNG 6 THÁNG
45+
ĐỐI TÁC TUYỂN DỤNG DOANH NGHIỆP
12\%
TIẾP TỤC HỌC CHƯƠNG TRÌNH TIẾN SĨ
Các Vị trí Tuyển dụng Hàng đầu \& Đối tác Doanh nghiệp Chính
Các Chức danh Công việc Mục tiêu Phổ biến
Senior Machine Learning Engineer (Kỹ sư Học máy Cao cấp)
AI Systems Architect (Kiến trúc sư Hệ thống AI)
Principal Data Scientist (Chuyên gia Khoa học Dữ liệu Chính)
MLOps Infrastructure Engineer (Kỹ sư Hạ tầng MLOps)
AI Regulatory \& Governance Specialist (Chuyên gia Quản trị \& Pháp lý AI)
Computer Vision / NLP Research Scientist (Nhà nghiên cứu Thị giác Máy tính / Xử lý Ngôn ngữ Tự nhiên)
Mạng lưới Đối tác Tuyển dụng
Google • Microsoft • Amazon Web Services • NVIDIA • Meta AI
• OpenAI • Snowflake • Databricks • Goldman Sachs •
JPMorgan Chase • Mayo Clinic Digital Health • Anthropic
22. Danh mục Kiểm tra Điều kiện Tốt nghiệp
Các Bước Hoàn thành Bằng cấp Cuối cùng:
1. Hoàn thành toàn bộ 36 tín chỉ bắt buộc với GPA tích lũy ≥ 3.00.
2. Nộp và bảo vệ thành công trước công chúng Luận văn / Đề án Tốt nghiệp (Capstone).
3. Hoàn tất xác nhận hoàn thành nghĩa vụ (clearance) từ Thư viện Trường và Cơ sở Siêu máy tính.
4. Nộp Đơn đăng ký Tốt nghiệp với Phòng Đào tạo (Office of the Registrar).
23. Thông tin Liên hệ Đơn vị
Để biết thêm thông tin chi tiết về chương trình học thuật, yêu cầu tuyển sinh hoặc hợp tác nghiên cứu, vui lòng liên hệ văn phòng chương trình:
Phòng Tuyển sinh Sau Đại học
Trường Máy tính, Viện Công nghệ Apex
100 Innovation Boulevard, Tech District
Email: admissions-dsai@apex.edu
Phòng Thường trực Học vụ
Điện thoại: +1 (555) 019-2831
Giờ làm việc: Thứ Hai – Thứ Sáu, 8:30 sáng – 5:00 chiều EST
Web: https://dsai.apex.edu/handbook
Thạc sĩ Khoa học Dữ liệu \& Trí tuệ Nhân tạo — Sổ tay Học thuật
Trang 12 / 12
\end{document}
